\subsection{模型的优点及创新}

\begin{enumerate}
    \item \textbf{区域差异化道路绕行系数。} 将全国城市按国家统计局标准划分为东部、中部、西部三大区域，分别设定1.25、1.40、1.60的道路绕行系数，相比单一绕行系数（如1.38）更准确地反映了不同地形条件下道路网络的实际弯绕特征。
    \item \textbf{往返空载成本的精确建模。} 区分去程（满载）与回程（空载）的油耗差异，设定空载油耗为满载的75\%，并结合路桥费、折旧、人工等固定成本的逐项核算，使运输成本计算更贴近实际运营。
    \item \textbf{工作天与日历天的区分。} 对于人工成本、轮胎损耗、折旧等与实际运营天数相关的固定费用项目，按年工作天数（280天）或维保摊消天数（300天）进行日成本分摊，区别于保险、审验等按日历天（365天）摊销的固定成本。
    \item \textbf{Clarke-Wright节约算法的方向约束增强。} 在经典CW节约算法基础上，引入方向偏离角约束、组内方向跨度约束和城市间距离上限约束，有效避免了"拼车但绕远路"的低效合并，使优化结果在实际路网中具有可操作性。
    \item \textbf{熵权法客观赋权的评估框架。} 采用基于信息熵的客观赋权方法，根据各指标在仓库间的实际差异程度自动分配权重，避免了主观打分可能带来的偏差；并通过TOPSIS辅助排名和Spearman秩相关检验验证了评估结果的稳健性。
    \item \textbf{分层评估与中间型指标设计。} 在评估指标中引入市内型与干线型配送的分层基准，并将配送密度设计为中间型指标——市内最优区间 $[12, 18]$ 点/日、干线最优区间 $[3, 5]$ 点/日，偏离区间两侧均线性扣分，避免了"越多越好"的片面评价。
\end{enumerate}

\subsection{关于模型的反思与改进}

\begin{enumerate}
    \item \textbf{道路距离模型的空间精度。} 当前使用Haversine球面距离乘以区域绕行系数的简化方法，未引入实际路网数据。未来可接入高德地图或百度地图API获取实际导航距离，进一步提高距离计算的精度。
    \item \textbf{装卸时间的简化假设。} 模型假设每次装卸固定为2小时，且忽略在多站路线中装卸顺序对时间窗口的影响。实际运营中，不同货物的装卸时间可能存在差异，严格的9:00–17:00窗口对到达时间有约束。
    \item \textbf{未考虑多仓库协同。} 问题二仅从南通单一仓库出发进行优化，而该企业在全国拥有5个枢纽仓库。在实际运作中，跨仓库协同配送和区域间库存调拨可以进一步提升效率。
    \item \textbf{成本参数的经验依赖性。} 模型中的空载油耗系数（0.75）、驱动日行驶上限（700\,km）等参数依赖运营经验，不同企业的实际情况可能有所不同，可通过历史运营数据的回归分析进行校准。
    \item \textbf{数据时间范围的局限性。} 问题三的评估数据仅涵盖一周（2018年12月10日至17日），季节性和周期性因素可能无法充分捕捉。更长时间跨度的数据有助于识别系统性规律和异常波动。
    \item \textbf{仅评估当前结构而不重新设计网络。} 问题三的优化在现有仓库布局和派车结构下进行，未从网络设计的角度考虑仓库选址、服务区域重新划分等更根本性的优化方向。
\end{enumerate}
